\documentclass[12pt]{article}
\usepackage{natbib}
\usepackage{floatrow,threeparttable}
\floatsetup[figure]{capposition=top}
\floatsetup[table]{capposition=top}
\usepackage{ragged2e}
\usepackage[left=1in, right=1in,top=1.25in,bottom=1.25in]{geometry}
\usepackage{multirow}
\usepackage{bbm}
\usepackage{lscape}
\usepackage{xparse}
\usepackage{xr}
\usepackage{amsmath,amssymb,amsthm,mathrsfs,amsfonts,dsfont}
\usepackage{commath}
\usepackage{graphicx}
\usepackage{booktabs}
\usepackage[bf]{titlesec}
\usepackage{setspace}
\usepackage{pdflscape}
\usepackage{bibunits}
\usepackage[nottoc,numbib]{tocbibind}
\defaultbibliography{bibliography.bib}
\defaultbibliographystyle{ecta}
\usepackage{verbatim}
\usepackage{subfigure}
\bibliographystyle{chicago}
\usepackage{soul}
\usepackage[addedmarkup=colored,deletedmarkup=sout]{changes}
\usepackage{tikz}
\usetikzlibrary{datavisualization}
\usetikzlibrary{datavisualization.formats.functions}
\usetikzlibrary{decorations.pathreplacing}
\tikzset{every picture/.style={line width=0.75pt}}
\usepackage{colortbl}
\usepackage[bf]{caption}
\usepackage{enumerate}
\usepackage[bookmarks, pdftitle={}, pdfauthor={}, colorlinks=true, linkcolor=red, citecolor=blue, urlcolor=blue]{hyperref}
\usepackage{txfonts}
\usetikzlibrary{decorations.pathreplacing}
\usepackage[multiple,stable]{footmisc}
\setlength{\footnotesep}{.1in}
\linespread{1.25}
\setlength{\parskip}{6pt}
\usepackage{datetime}
\newdateformat{monthyeardate}{\monthname[\THEMONTH], \THEYEAR}

\begin{document}

%%%%%%%%%%%%%%%%%%%%%%%%%%%%%%%%%%%%%%%%%
%%                TITLE                %%
%%%%%%%%%%%%%%%%%%%%%%%%%%%%%%%%%%%%%%%%%

\title{\Large{Do CCT Programs (Really) Reduce Mortality? Ten-Year Evidence from Mexico}%
\thanks{This version: March 2026. The opinions expressed in this paper are those of the authors and do not necessarily reflect the views of the Inter-American Development Bank, its Board of Directors, or the countries they represent, nor of INEGI. All remaining errors are our own. All data used in this study are publicly available from INEGI.}}

\author{Emma Aguila\footnote{Sol Price School of Public Policy, University of Southern California.}
\and William H.\ Dow\footnote{School of Public Health and Department of Demography, University of California, Berkeley.}
\and Felipe Menares\footnote{Inter-American Development Bank. Email: \texttt{felipeme@iadb.org}}
\and Susan W.\ Parker\footnote{School of Public Policy, University of Maryland.}
\and Jorge Peniche\footnote{Sol Price School of Public Policy, University of Southern California.}
\and Soomin Ryu\footnote{Department of Public Health, University of Missouri.}}

\date{}
\maketitle
\vspace{-12mm}
\thispagestyle{empty}

%%%%%%%%%%%%%%%%%%%%%%%%%%%%%%%%%%%%%%%%%
%%               ABSTRACT              %%
%%%%%%%%%%%%%%%%%%%%%%%%%%%%%%%%%%%%%%%%%

\begin{abstract}
\noindent
Early evidence on Mexico's Progresa conditional cash transfer program suggested large short-term reductions in elderly mortality. We re-examine these effects over ten years using difference-in-differences models that exploit geographic variation in program rollout across highly marginalized municipalities, combining twenty years of vital statistics with administrative records on program enrollment. We find no evidence of significant effects of Progresa on the elder mortality rate (EMR) among adults aged 65 and older. Replicating the specification of \cite{BarhamRowberry2013JDE}, we show their results are sensitive to weighting and lag structure, becoming smaller and statistically insignificant with minor changes. These null results are arguably not surprising: program benefits were primarily conditioned on children's schooling, the elderly were not direct recipients of transfers in most cases, and health conditionalities required only one annual clinic visit for adults.

\noindent \\
\noindent\textbf{Keywords:} Mortality, Conditional Cash Transfers, Progresa, Mexico, Elderly \\
\noindent\textbf{JEL Codes:} I13, I14, I15, I18
\end{abstract}

\clearpage
\setcounter{page}{1}

%%%%%%%%%%%%%%%%%%%%%%%%%%%%%%%%%%%%%%%%%
%%            INTRODUCTION             %%
%%%%%%%%%%%%%%%%%%%%%%%%%%%%%%%%%%%%%%%%%

\section{Introduction}\label{s:intro}

The health effects of income-support programs for older adults have attracted substantial attention, yet the evidence remains mixed. Some studies find important health benefits \citep{Aguila2015PNAS, Case2004PerspectivesAging}, others find no effects \citep{LloydSherlockAgrawal2014JDS}, and still others identify unintended adverse consequences linked to obesity and diet change \citep{FernaldGertlerHou2008JN, KronebuschDamon2019EHB}. Most existing evidence focuses on short-run outcomes, leaving long-run mortality effects poorly understood.

Mexico's Progresa conditional cash transfer (CCT) program, launched in 1997 and subsequently renamed Oportunidades and then Prospera, offers a rare opportunity to study longer-term mortality effects of a large-scale transfer program. Progresa has served as a model for CCT programs in over 60 countries \citep{Baird2014JDEffectiveness, Fiszbein2009CCTBook, ParkerTodd2017JEL}, and a growing literature evaluates its health impacts \citep{Gertler2004AER, LagardeHainesPalmer2009Cochrane, ParkerTodd2017JEL, RiveraHernandez2016HSR}. However, long-term effects—particularly on elderly mortality—remain largely unstudied \citep{ParkerTodd2017JEL}.

We are aware of only a small number of studies examining health and mortality effects among the aging under a CCT program. \cite{BarhamRowberry2013JDE} found that two years after implementation, Progresa reduced municipal-level mortality rates for adults over 65 by an average of 4\%, primarily through decreases in infectious disease and diabetes-related deaths. \cite{BehrmanParker2013Demography} found that health utilization outcomes improved for elderly women but not men six years after implementation. Using the same Progresa program, \cite{Aguila2024EHB} find that the program improved hypertension diagnosis and treatment use among adults aged 50 and older but did not reduce uncontrolled blood pressure as measured by biomarkers—suggesting that while CCTs can raise healthcare utilization, they may fall short of producing clinically meaningful improvements in chronic disease control. In contrast, related evidence on Mexico's subsequent non-contributory pension program \textit{70 y Más}, which provided direct transfers to adults aged 70 and older, finds significant reductions in elderly mortality in rural areas \citep{Menares2026JHR}. The contrast between a pension targeting the elderly directly and a CCT structured around children's conditionalities motivates our re-examination of Progresa's mortality effects over a longer horizon.

In this paper, we estimate difference-in-differences models exploiting variation in the proportion of households enrolled in Progresa across municipalities, comparing municipalities that enrolled more intensively in the program's early phase (1997--1999) to those that enrolled more intensively in the later phase (2000--2005), following the strategy of \cite{parker2023conditional}. We focus on highly marginalized municipalities where the program was most heavily targeted, and study the period 1991--2006, stopping before a major structural change in 2006 when Progresa added a direct pension component for the elderly.

Our main finding is that Progresa had no significant effect on elderly mortality over its first ten years of operation. This result is robust to weighting, inclusion of Seguro Popular controls, and alternative outcome measures. We further replicate \cite{BarhamRowberry2013JDE} and show that their positive estimates are sensitive to specification choices. We argue these null results are consistent with program design: benefits were primarily conditioned on children's schooling, elderly individuals living in extended households were not the direct transfer recipients, and health conditionalities required only one annual clinic visit.

The paper proceeds as follows. Section~\ref{s:progresa} describes the Progresa program. Section~\ref{s:data} presents the data and empirical strategy. Section~\ref{s:results} reports the main results and robustness checks. Section~\ref{s:conclusion} concludes.

%%%%%%%%%%%%%%%%%%%%%%%%%%%%%%%%%%%%%%%%%
%%         PROGRAM BACKGROUND          %%
%%%%%%%%%%%%%%%%%%%%%%%%%%%%%%%%%%%%%%%%%

\section{Mexico's Progresa CCT Program}\label{s:progresa}

Progresa began in small rural communities in 1997, providing monetary transfers to poor families conditional on children attending school and family members visiting health clinics. The program rapidly expanded to cover six million families—one quarter of all Mexican families—with two thirds of beneficiaries residing in communities with fewer than 2,500 inhabitants. Average monthly transfers reached approximately MXN~\$800 (US~\$90 PPP), representing more than a 25\% income increase for many poor households. All transfers are paid to the female head of household.

The program includes two main grant types: grants conditional on children's school attendance (for children under 22), and a health and nutrition transfer requiring all family members to attend regular health clinic visits. For adults aged 17 and older, including the elderly, the attendance requirement is one clinic visit per year—a minimal conditionality. In 2006, a pension component for adults aged 70 and older was added, fundamentally altering the program's reach for the elderly. We end our analysis in 2006 to focus on the original CCT design.

The majority of elderly individuals in rural eligible households live in extended families \citep{BehrmanParker2013Demography}. Before the program, approximately 75\% of elderly individuals lived in households with children under 15. In most cases, elderly individuals are grandparents of children receiving educational grants and thus do not directly control program expenditures, though they may benefit indirectly. Given this structure, large and immediate effects on elderly health and mortality would not be expected.

Elderly mortality in Mexico has declined over time and is slightly higher in non-marginalized areas than in marginalized (poor) areas, as shown in Figure~\ref{f:variation}. \cite{BeltranSanchezWong2025CSP} demonstrate using panel data from the Mexican Health and Aging Survey that this pattern is not explained by selective migration of unhealthy individuals toward urban areas.

%%%%%%%%%%%%%%%%%%%%%%%%%%%%%%%%%%%%%%%%%
%%          DATA AND METHODS           %%
%%%%%%%%%%%%%%%%%%%%%%%%%%%%%%%%%%%%%%%%%

\section{Data and Empirical Strategy}\label{s:data}

\subsection{Data}\label{s:sources}

\textbf{Mortality data.} We use individual-level death registries derived from death certificates, covering every death in Mexico between 1991 and 2006—over three million records. The data include cause of death, birth year, sex, and the municipality of residence of the deceased, which we use to link to population denominators and program variables. All data are publicly available from the Mexican National Institute of Statistics and Geography (INEGI).\footnote{INEGI derives mortality data from the certification system of the Mexican Ministry of Public Health. Mexico generally maintains high-quality vital statistics \citep{mahapatra2007civil, mikkelsen2015global}.}

The expansion of Seguro Popular—Mexico's public health insurance program for informal-sector workers, which rolled out staggered across municipalities beginning in 2002—may have improved the accuracy of death registration. To address this, we control for Seguro Popular enrollment intensity at the municipality level.\footnote{Seguro Popular reached full municipal coverage by 2006 but significant portions of the eligible population remained uninsured throughout our study period.}

\textbf{Population data.} We use INEGI population counts disaggregated by sex, municipality, and age from Census data (1990, 2000, 2005, 2010, 2020), linearly interpolating annual estimates for each sex-age-municipality cell. We similarly interpolate municipal household counts to construct program intensity denominators.

\textbf{Progresa administrative data.} We obtain annual counts of enrolled Progresa households by municipality from administrative records, which we use to construct our treatment intensity variable.

\textbf{Additional controls.} We use the 1990 CONAPO Municipal Marginalization Index to restrict the sample to highly marginalized areas and as a covariate. We also incorporate administrative data on Seguro Popular beneficiaries to construct a municipality-year enrollment intensity variable.\footnote{The Marginalization Index classifies all Mexican municipalities into five categories based on household poverty indicators, including dirt floors, crowding, illiteracy, and lack of sanitation infrastructure.}

\subsection{Sample and Descriptive Statistics}\label{s:sample}

Our analysis sample is restricted to deaths among adults aged 65 and older, in highly marginalized municipalities classified as high (grade 4) or very high (grade 5) on the 1990 CONAPO Marginalization Index. We construct municipality-year cells combining death counts with population denominators and program variables. Table~\ref{t:descriptives} reports summary statistics comparing marginalized and non-marginalized municipalities in the pre-program period. As expected, marginalized municipalities have substantially worse socioeconomic conditions and much higher Progresa enrollment rates.

\subsection{Empirical Strategy}\label{subsec:empirical_strategy}

We exploit variation in the proportion of households enrolled in Progresa across municipalities before and after the program began in 1997. Following \cite{parker2023conditional}, who study the long-run effects of Progresa on children's educational and labor market outcomes, we distinguish between municipalities that enrolled households more intensively in the early phase (1997--1999, denoted $\textit{Intensity}_{1999}$) and those that enrolled more intensively in the later phase (2000--2005, denoted $\textit{Intensity}_{2005}$). Parker and Vogl argue that after 2005 program expansion shifted substantially toward urban areas, making the two phases distinct and largely non-overlapping sources of intensity variation across municipalities. We adopt the same periodization and apply it to elderly mortality.

Because Progresa was preferentially targeted toward more disadvantaged areas, a simple difference-in-differences comparing Progresa municipalities to all others could be confounded by differential trends. Restricting the sample to highly marginalized municipalities mitigates this concern, as identification then requires parallel trends only between early and late roll-out municipalities within this group. Our estimates capture the impact of earlier versus later cumulative exposure to the program, under the assumption that additional years of access accumulate into mortality effects. We estimate:

\begin{equation}\label{eq:did}
\text{EMR}_{m,t}
= \beta_{0}\,\textit{Post}_{t}\times\textit{Intensity}_{m,1999}
+ \beta_{1}\,\textit{Post}_{t}\times\textit{Intensity}_{m,2005}
+ \delta_{m} + \gamma_{t} + \chi_{m,t} + \varepsilon_{m,t}
\end{equation}

\noindent where $\text{EMR}_{m,t}$ is the Elder Mortality Rate (deaths per 1,000 inhabitants aged 65+) in municipality $m$ in year $t$. $\textit{Post}_{t}$ equals one for years 1997--2006 and zero for 1991--1996. Municipality fixed effects $\delta_{m}$ control for time-invariant unobservables; year fixed effects $\gamma_{t}$ absorb common time trends. $\chi_{m,t}$ is the Seguro Popular enrollment intensity, included in specifications that control for this potentially confounding program. Standard errors are clustered at the municipality level.

We also estimate event study specifications to trace dynamic effects and test for differential pre-trends:

\begin{equation}\label{eq:eventstudy}
\text{EMR}_{m,t}
= \sum_{k \neq 1996} \beta_{k}\,\mathbf{1}[t=k]\times\textit{Intensity}_{m,1999}
+ \sum_{k \neq 1996} \theta_{k}\,\mathbf{1}[t=k]\times\textit{Intensity}_{m,2005}
+ \delta_{m} + \gamma_{t} + \chi_{m,t} + \varepsilon_{m,t}
\end{equation}

\noindent with 1996 as the reference year. The $\beta_{k}$ coefficients trace the effect of early-phase enrollment intensity on mortality in each year relative to 1996.

%%%%%%%%%%%%%%%%%%%%%%%%%%%%%%%%%%%%%%%%%
%%              RESULTS                %%
%%%%%%%%%%%%%%%%%%%%%%%%%%%%%%%%%%%%%%%%%

\section{Results}\label{s:results}

\subsection{Main Results}\label{s:main_results}

Table~\ref{t:did} reports difference-in-differences estimates of the effect of Progresa enrollment intensity on the EMR for adults aged 65 and older. Panel~A reports pooled results, Panel~B results for males, and Panel~C results for females. Within each panel, column~(1) presents unweighted estimates without Seguro Popular controls; column~(2) adds Seguro Popular intensity; column~(3) presents population-weighted estimates (weighted by the elderly population); and column~(4) combines weighting with Seguro Popular controls.

Across all specifications, columns, and panels, the estimated effect of Progresa enrollment intensity on elderly mortality is small and statistically insignificant. The interaction of early-phase intensity ($\textit{Intensity}_{1999} \times \textit{Post}$) is close to zero in both the unweighted and weighted specifications, for the pooled sample and separately by sex. Adding Seguro Popular controls in columns~(2) and~(4) does not materially alter results. These findings hold regardless of whether we weight by the size of the elderly population.

Figure~\ref{f:es_dd_sex_no_weight} presents the event study estimates for the weighted specification with Seguro Popular controls (column~4). The coefficients are imprecise but reveal no systematic pattern: pre-program coefficients do not exhibit statistically significant trends, supporting the parallel trends assumption. Post-program coefficients show no consistent negative effect, and the pattern is similar for males and females. The null result is therefore not an artifact of the choice of reference year or the post-period definition.

\subsection{Replication and Robustness}\label{s:robustness}

Our results contrast with \cite{BarhamRowberry2013JDE}, who found that a one percentage point increase in two-year lagged Progresa intensity reduced elderly mortality by approximately 0.064 deaths per 1,000. To explore this discrepancy, we replicate their specification:

\begin{equation}\label{eq:br}
\text{EMR}_{m,t}
= \alpha_{0} + \alpha_{1}\,\textit{Intensity}_{m,t-2}
+ \delta_{m} + \gamma_{t} + \chi_{m,t} + \varepsilon_{m,t}
\end{equation}

Appendix Table~\ref{t:did_robust_br} replicates their exercise restricting to municipalities incorporated into Progresa in 1998 or 1999. Appendix Table~\ref{t:did_robust_br2} extends the sensitivity analysis across alternative lag structures, weighting schemes, and sample restrictions. Row~(A) reproduces \cite{BarhamRowberry2013JDE}'s original estimates; row~(B) shows our replication, which yields similar but not identical results, with a negative and significant coefficient of approximately $-0.07$. Adding population weights attenuates the estimate by roughly 50\%. Using one-year or three-year lags instead of the two-year lag further reduces the estimate to approximately one third of the original magnitude, and precision falls substantially. These results suggest that \cite{BarhamRowberry2013JDE}'s finding is sensitive to relatively minor specification choices.

Appendix Figure~\ref{f:es_cod_marg_a}, \ref{f:es_cod_marg_b}, \ref{f:es_cod_br_a}, and~\ref{f:es_cod_br_b} examine cause-specific mortality—cancer, cardiovascular disease, respiratory disease, infectious disease, diabetes, nutritional causes, ill-defined causes, and accidents—separately for the highly marginalized sample and the Barham-Rowberry sample. Appendix Figures~\ref{f:es_dd_sex_no_weight_app2} and~\ref{f:es_dd_sex_no_weight_app3} present additional event study specifications using alternative mortality measures (Age-Adjusted Mortality Rate, AAMR) and samples. Across these robustness checks, we consistently find no evidence of significant program effects on elderly mortality.

\subsection{Mechanisms: Labor Supply and Household Composition}\label{s:mechanisms}

The null mortality finding raises the question of whether Progresa moved any of the intermediate variables that could plausibly link income transfers to elderly health. Two candidate channels are particularly relevant for this age group: relief from labor burden and changes in household composition. If the program reduced hours worked by elderly members—freeing time for rest, preventive care, and healthcare-seeking—or if it encouraged adult children to remain in or return to the household, providing informal caregiving, these shifts could translate into lower mortality. On the other hand, if neither channel was activated, the null mortality result is more readily explained.

Table~\ref{tab:combined_elderly} tests these pathways using experimental data from eligible households in the original 1997 Progresa rollout, estimating difference-in-differences models with 1997 as the baseline. Column~(1) reports the effect on weekly hours worked by elderly individuals; columns~(2)--(4) report effects on the probability of living alone, living with children, and living exclusively with other elderly members, respectively. The estimates are broadly consistent with the main mortality finding: Progresa did not generate systematic changes in labor supply or household composition among elderly members of eligible households. The program did not significantly alter whether elderly individuals co-resided with children, the living arrangement most likely to provide informal caregiving and social support, and the absence of this effect reinforces why mortality impacts would be limited \citep{BehrmanParker2013Demography, duflo2000child}. These results complement the evidence in \cite{Aguila2024EHB} that Progresa improved healthcare access but not downstream health biomarkers, suggesting the program operated near the margins of what a child-focused CCT can do for an older adult population.

%%%%%%%%%%%%%%%%%%%%%%%%%%%%%%%%%%%%%%%%%
%%             CONCLUSION              %%
%%%%%%%%%%%%%%%%%%%%%%%%%%%%%%%%%%%%%%%%%

\section{Conclusion}\label{s:conclusion}

This paper re-examines the medium-term effect of Mexico's Progresa conditional cash transfer program on elderly mortality. Using twenty years of vital statistics and administrative enrollment records, we find no evidence that Progresa significantly reduced mortality among adults aged 65 and older during its first ten years of operation.

These null results are plausible given program design. Most transfers in Progresa are linked to children's school attendance, effectively subsidizing investments in children rather than directly supporting the elderly. Because elderly individuals in eligible households typically live in extended families, they rarely control program resources. The health conditionality for aging adults requires only one annual clinic visit, a minimal intervention. Consistent with this interpretation, Table~\ref{tab:combined_elderly} shows that Progresa did not materially alter elderly labor supply or household composition—the two intermediate channels most likely to translate transfer income into mortality improvements. This absence of behavioral response, combined with the evidence in \cite{Aguila2024EHB} that Progresa improved hypertension treatment access but not actual blood pressure control, paints a coherent picture: a child-focused CCT can raise health care utilization at the margin but is unlikely to generate the sustained changes needed to reduce elderly mortality.

We further show that the key prior finding of \cite{BarhamRowberry2013JDE} is sensitive to specification. Their result attenuates by 50\% with population weighting and shrinks further with alternative lag structures, losing statistical significance. Combined with our own estimates, we conclude that the evidence for large mortality effects of conditional cash transfers on older adults is not strong. In contrast, \cite{Menares2026JHR} find that Mexico's \textit{70 y Más} non-contributory pension—which transferred resources directly and unconditionally to adults aged 70 and older—did reduce elderly mortality in rural areas. Taken together, the evidence suggests that program design matters: transfers conditioned on children's schooling and targeted indirectly to the elderly are unlikely to generate mortality benefits comparable to pensions that place income directly in the hands of older adults \citep{GalianiGertlerBando2016LE, Aguila2015PNAS}.

\newpage

%%%%%%%%%%%%%%%%%%%%%%%%%%%%%%%%%%%%%%%%%
%%             REFERENCES              %%
%%%%%%%%%%%%%%%%%%%%%%%%%%%%%%%%%%%%%%%%%

\setcounter{secnumdepth}{0}
\bigskip
\singlespace
\renewcommand{\bibname}{References}
{\linespread{1.0}\selectfont\bibliography{bibliography}}

%%%%%%%%%%%%%%%%%%%%%%%%%%%%%%%%%%%%%%%%%
%%              FIGURES                %%
%%%%%%%%%%%%%%%%%%%%%%%%%%%%%%%%%%%%%%%%%

\newpage
%=======================================================================
% figures.tex — Pre-Trend Event Study Figures (Mechanisms, ENIGH)
% Script: 05_mechanisms_eventstudy_enigh.do
% Layout: 2 columns x 4 rows per panel (max 8 subfigures)
% Paths:  pretrend/{marg,all}/F_PT_{1998,2000}_{outcome}
% Requires in main preamble: \usepackage{graphicx,subfigure}
%=======================================================================

% Compact subfigure helper: \ptsfig{filepath}{subcaption}
\newcommand{\ptsfig}[2]{%
  \subfigure[#2]{\includegraphics[width=0.47\linewidth]{#1}}}


%-----------------------------------------------------------------------


\begin{figure}[H]
    \caption[]{The effects of interaction terms between Progresa in 1999 and years on elderly mortality +65 overall and by sex in marginalized municipalities, Mexico, 1991-2006. }
    \label{f:es_dd_sex_no_weight}

  \subfigure[All]
{\includegraphics[width=0.69\textwidth]{figures/main/Figure_2a_all.pdf}}
    
    \subfigure[Males]
     {\includegraphics[width=0.49\textwidth]{figures/main/Figure_2b_male.pdf}}
    \subfigure[Females]
     {\includegraphics[width=0.49\textwidth]{figures/main/Figure_2c_female.pdf}}
     
         \floatfoot{Notes: The data analysis is restricted to only marginalized municipalities, 1991-2006. The plots present the interaction terms between Progresa intensity in 1999 and years. Year and municipality fixed effects are included. Standard errors are clustered at the municipality level. 1996 is a reference year and the regression models are not weighted.
Source: Authors’ calculations using data on deaths and population from the Instituto Nacional de Estadistica y Geografia and Mexican Census data.
}
\end{figure}

\begin{figure}[H]
    \caption[]{The effects of interaction terms between Progresa in 1999 and years on elderly mortality +65 overall and by sex in marginalized municipalities, Mexico, 1991-2006. }
    \label{f:es_dd_sex_no_weight}

  \subfigure[All]
{\includegraphics[width=0.69\textwidth]{figures/main/Figure_3a_all.pdf}}
    
    \subfigure[Males]
     {\includegraphics[width=0.49\textwidth]{figures/main/Figure_3b_male.pdf}}
    \subfigure[Females]
     {\includegraphics[width=0.49\textwidth]{figures/main/Figure_3c_female.pdf}}
     
         \floatfoot{Notes: The data analysis is restricted to only marginalized municipalities, 1991-2006. The plots present the interaction terms between Progresa intensity in 1999 and years. Year and municipality fixed effects are included. Standard errors are clustered at the municipality level. 1996 is a reference year and the regression models are not weighted.
Source: Authors’ calculations using data on deaths and population from the Instituto Nacional de Estadistica y Geografia and Mexican Census data.
}
\end{figure}


\begin{figure}[htbp]\centering
\ptsfig{figures/pretrend/all/F_PT_1998_employed}{Employment}\hfill
\ptsfig{figures/pretrend/all/F_PT_1998_hrs_worked}{Hrs Worked}\\[1ex]
\ptsfig{figures/pretrend/all/F_PT_1998_hrs_worked_pos}{Hrs Worked ($>$0)}\hfill
\ptsfig{figures/pretrend/all/F_PT_1998_ind_earnings}{Earnings}\\[1ex]
\ptsfig{figures/pretrend/all/F_PT_1998_ind_income_tot}{Income}\hfill
\ptsfig{figures/pretrend/all/F_PT_1998_progresa_ind}{PROGRESA Transfers}
\caption{Even Study: Individual Labor Outcomes - Intensity 1998, All Municipalities}
\label{fig:pt_1998_all_t1}
\end{figure}

\begin{figure}[htbp]\centering
\ptsfig{figures/pretrend/all/F_PT_1998_hh_earnings}{Earnings}\hfill
\ptsfig{figures/pretrend/all/F_PT_1998_hh_income_tot}{Income}\\[1ex]
\ptsfig{figures/pretrend/all/F_PT_1998_hh_expenditure}{Expenditure}\hfill
\ptsfig{figures/pretrend/all/F_PT_1998_progresa_hh}{PROGRESA}\\[1ex]
\ptsfig{figures/pretrend/all/F_PT_1998_savings}{Savings}\hfill
\ptsfig{figures/pretrend/all/F_PT_1998_debt}{Debt}
\caption{Even Study: Household Labor Outcomes - Intensity 1998, All Municipalities}
\label{fig:pt_1998_all_t2}
\end{figure}

\begin{figure}[htbp]\centering
\ptsfig{figures/pretrend/all/F_PT_1998_food_exp}{Total Food}\hfill
\ptsfig{figures/pretrend/all/F_PT_1998_vegg_fruit}{Veggies \& Fruit}\\[1ex]
\ptsfig{figures/pretrend/all/F_PT_1998_cereals}{Cereals}\hfill
\ptsfig{figures/pretrend/all/F_PT_1998_meat_dairy}{Meat \& Dairy}\\[1ex]
\ptsfig{figures/pretrend/all/F_PT_1998_sugar_fat_drink}{Sugar \& Fat}\hfill
\ptsfig{figures/pretrend/all/F_PT_1998_vice}{Alcohol and Tobacco}
\caption{Event Study: Household Food Outcomes - Intensity 1998, All Municipalities}
\label{fig:pt_1998_all_t3}
\end{figure}

\begin{figure}[htbp]\centering
\ptsfig{figures/pretrend/all/F_PT_1998_health_exp}{Total Health}\hfill
\ptsfig{figures/pretrend/all/F_PT_1998_medical}{Medical Visits}\\[1ex]
\ptsfig{figures/pretrend/all/F_PT_1998_medical_inpatient}{Inpatient}\hfill
\ptsfig{figures/pretrend/all/F_PT_1998_medical_outpatient}{Outpatient}\\[1ex]
\ptsfig{figures/pretrend/all/F_PT_1998_drugs}{Drugs (Rx and OTC)}\hfill
\ptsfig{figures/pretrend/all/F_PT_1998_ortho}{Orthotics}
\caption{Event Study: Household Health Outcomes - Intensity 1998, All Municipalities.}
\label{fig:pt_1998_all_t4}
\end{figure}






\begin{figure}[H]
    \caption[]{Share of progresa penetration}
    \label{f:es_dd_sex_no_weight}

    \subfigure[1997 - All]
     {\includegraphics[width=0.49\textwidth]{figures/appendix/FA0a_inten1997_kdensity.pdf}}
    \subfigure[1997 - by wave]
     {\includegraphics[width=0.49\textwidth]{figures/appendix/FA0b_inten1997_bywave.pdf}}
    

    \subfigure[1999 - All]
     {\includegraphics[width=0.49\textwidth]{figures/appendix/FA0c_inten1999_kdensity.pdf}}
    \subfigure[1999 - by wave]
     {\includegraphics[width=0.49\textwidth]{figures/appendix/FA0d_inten1999_bywave.pdf}}

    \subfigure[Continous - by Wave ENIGH Sample]
     {\includegraphics[width=0.49\textwidth]{figures/appendix/FA0f_intensity_new_surveytrend}}
    \subfigure[Continious - Mortality Sample]
     {\includegraphics[width=0.49\textwidth]{figures/appendix/FA0h_intensity_new_annual_trend}}
     
         \floatfoot{Notes: }
\end{figure}


%%%%%%%%%%%%%%%%%%%%%%%%%%%%%%%%%%%%%%%%%
%%              TABLES                 %%
%%%%%%%%%%%%%%%%%%%%%%%%%%%%%%%%%%%%%%%%%

\newpage

\begin{table}[H]
\RawCaption%
  {\caption[]{Summary statistics of non-marginalized and marginalized municipalities, Mexico (\%)}
   \label{t:descriptives}}
\floatbox[{\captop}]{table}[\FBwidth]
  {}
  {\scalebox{0.9}{
      \begin{tabular}{lcc}
\hline
 & \textbf{Non-marginalized areas} & \textbf{Marginalized areas} \\
\hline

\multicolumn{3}{l}{\textbf{A. Progresa program intensity}} \\
\quad 1999 & 0.10 (0.13) & 0.44 (0.22) \\
\quad 2005 & 0.28 (0.19) & 0.72 (0.17) \\[6pt]

\multicolumn{3}{l}{\textbf{B. Characteristics of municipalities, 1990}} \\
\multicolumn{3}{l}{\quad \textit{Components of the marginality index}} \\
\quad Illiterate population & 13.33 (6.57) & 33.33 (13.3) \\
\quad With less than primary education & 45.74 (12.32) & 69.56 (9.68) \\
\quad Without a toilet & 25.94 (16.62) & 60.30 (18.54) \\
\quad Without electricity & 11.81 (10.29) & 36.39 (24.72) \\
\quad Without running water & 19.43 (14.87) & 50.65 (24.07) \\
\quad With crowding & 59.84 (9.79) & 73.96 (8.24) \\
\quad With dirt floor & 22.05 (14.21) & 61.98 (21.58) \\
\quad In communities with $<$5{,}000 inhabitants & 60.18 (36.15) & 95.50 (13.54) \\
\quad Earning $<$2× minimum wage & 69.19 (11.60) & 85.82 (8.55) \\[6pt]

Number of municipalities & 1,234 & 1,119 \\
\hline
\end{tabular}

    }
    \vspace{-3mm}
    \floatfoot{Notes: Section A reports the mean proportion (\%) with standard deviations of municipalities with a ratio of the Progresa beneficiary households to the total number of households in the municipality. Section B presents sample mean (\%) of each variable and standard deviations are in parentheses. Marginalized areas include municipalities categorized as high and very high level of the Margination Index, while non-marginalized areas include municipalities categorized as medium, low, and very low level of the Margination Index. “With crowding” is measured by number of rooms divided by household size. Source: Authors’ calculations using Mexican Census 1990 and Progresa administrative data on beneficiaries.}
  }
\end{table}


\begin{table}[H]
\RawCaption%
  {\caption[]{Effects of Progresa on Elderly Mortality (65+) by Sex, Marginalized Municipalities}
   \label{t:did}}
\floatbox[{\captop}]{table}[\FBwidth]
  {}
  {\scalebox{0.9}{
      \begin{tabular}{lcccc} \hline \hline
& \multicolumn{1}{c}{(1)} & \multicolumn{1}{c}{(2)} & \multicolumn{1}{c}{(3)} & \multicolumn{1}{c}{(4)} \\ \toprule
\underline{\textit{Panel A: Pooled}}  \\ 
\textit{Intensity 1999 x post (1997-2006)} &        0.369 &        0.354 &       -0.777 &       -0.719 \\ 
 & (       2.546) & (       2.547) & (       1.464) & (       1.464) \\ 
  & & & & \\ 
\textit{Intensity 2005 x post (1997-2006)} &       -5.525* &       -5.513* &       -2.787 &       -2.796 \\ 
 & (       3.075) & (       3.074) & (       1.806) & (       1.803) \\ 
  & & & & \\ 
Mean (1991-1996) &        46.47 &        46.47 &        46.47 &        46.47 \\ 
Obs &       17,904 &       17,904 &       17,904 &       17,904 \\ 
  & & & & \\ 
\underline{\textit{Panel B: Males}}  \\ 
\textit{Intensity 1999 x post (1997-2006)} &       -1.627 &       -1.588 &       -2.288 &       -2.206 \\ 
 & (       2.767) & (       2.776) & (       1.732) & (       1.735) \\ 
  & & & & \\ 
\textit{Intensity 2005 x post (1997-2006)} &       -1.488 &       -1.519 &        0.214 &        0.201 \\ 
 & (       3.309) & (       3.312) & (       2.000) & (       1.994) \\ 
  & & & & \\ 
Mean (1991-1996) &        47.26 &        47.26 &        47.26 &        47.26 \\ 
Obs &       17,904 &       17,904 &       17,904 &       17,904 \\ 
  & & & & \\ 
\underline{\textit{Panel C: Females}}  \\ 
\textit{Intensity 1999 x post (1997-2006)} &        1.737 &        1.668 &        0.592 &        0.620 \\ 
 & (       2.932) & (       2.928) & (       1.824) & (       1.824) \\ 
  & & & & \\ 
\textit{Intensity 2005 x post (1997-2006)} &       -9.117** &       -9.063** &       -5.637** &       -5.641** \\ 
 & (       3.776) & (       3.770) & (       2.205) & (       2.205) \\ 
  & & & & \\ 
Mean (1991-1996) &        46.11 &        46.11 &        46.11 &        46.11 \\ 
Obs &       17,904 &       17,904 &       17,904 &       17,904 \\ 
  & & & & \\ 
No. Mun &        1,119 &        1,119 &        1,119 &        1,119 \\ 
  & & & & \\ 
Year FE & Y & Y & Y & Y \\ 
Mun FE & Y & Y & Y & Y \\ 
Seguro Popular & N & Y & N & Y \\ 
Weights & N & N & Y & Y \\ 
Cluster SE: Mun & Y & Y & Y & Y \\ 
\bottomrule
\end{tabular}
    }
    \vspace{-3mm}
    \floatfoot{Notes: The analysis is restricted to marginalized municipalities, 1991–2006. Progresa and Seguro Popular intensities are continuous variables. The post-dummy equals 1 for years 1997–2006. Standard errors are clustered at the municipality level and shown in parentheses. ***p$<0.01$, **p$<0.05$, *p$<0.1$.
}
  }
\end{table}


\begin{table}[htbp]
\centering
\caption{Progresa Impacts on Labor Supply and Living Arrangements: Eligible Elderly (Age 65+ in 1997)}
\label{tab:combined_elderly}
\begin{threeparttable}

\begin{tabular}{lcccc}
\toprule
 & (1) & (2) & (3) & (4) \\
 & Weekly Hours & Live Alone & With Children & Only Elderly \\
\midrule
\multicolumn{5}{l}{\textbf{Panel A: Females}} \\
\midrule
Treat $\times$ 1998 & 0.124 & 0.000 & 0.001 & $-$0.008 \\
                    & (1.136) & (0.008) & (0.011) & (0.008) \\[4pt]
Treat $\times$ 1999 & 0.292 & 0.007 & 0.015 & 0.012 \\
                    & (1.070) & (0.010) & (0.018) & (0.014) \\[4pt]
Observations        & 3,338 & 3,597 & 3,597 & 3,597 \\
Control Mean (1997) & 3.923 & 0.044 & 0.708 & 0.145 \\
\midrule
\multicolumn{5}{l}{\textbf{Panel B: Males}} \\
\midrule
Treat $\times$ 1998 & $-$1.210 & 0.012 & $-$0.008 & $-$0.007 \\
                    & (1.895) & (0.008) & (0.011) & (0.010) \\[4pt]
Treat $\times$ 1999 & 0.378 & 0.008 & 0.040$^{**}$ & $-$0.001 \\
                    & (1.872) & (0.012) & (0.020) & (0.017) \\[4pt]
Observations        & 3,170 & 3,419 & 3,419 & 3,419 \\
Control Mean (1997) & 25.288 & 0.049 & 0.676 & 0.159 \\
\bottomrule
\end{tabular}

\begin{tablenotes}
\small
\item \textit{Notes:} Difference-in-differences estimates with 1997 as baseline. Sample restricted to 
eligible households. Columns (1) reports weekly hours worked. Columns (2)--(4) report living 
arrangement outcomes: living alone, living with children, and living only with other elderly, 
respectively. Standard errors in parentheses.
\item $^{*}$ $p < 0.10$, $^{**}$ $p < 0.05$, $^{***}$ $p < 0.01$
\end{tablenotes}
\end{threeparttable}

\end{table}

\appendix

\newpage
\setcounter{page}{1}
\section*{Online Appendix}\label{appOnline}
\vspace{2cm}
\begin{center}
{\Large \it Do CCT Programs (Really) Reduce Mortality: Ten-Year Evidence from Mexico} \\[.4cm]
Emma Aguila, William H.\ Dow, Felipe Menares, Susan Parker, Jorge Peniche, and Soomin Ryu
\end{center}
\listoffigures
\listoftables

\newpage
\section*{Appendix A: Additional Figures and Tables}\label{appA}
\setcounter{table}{0}
\renewcommand{\thetable}{A.\arabic{table}}
\setcounter{figure}{0}
\renewcommand{\thefigure}{A.\arabic{figure}}
\addtocontents{toc}{\setcounter{tocdepth}{1}}
\renewcommand{\thesubsection}{\Alph{subsection}}



\begin{figure}[H]
    \caption{Mortality and Progresa Intensity in Mexico}
    \label{f:variation}

    \subfigure[Mortality and Progresa Intensity]
    {\includegraphics[width=0.69\textwidth]{figures/appendix/Figure_1a_all.pdf}}

    \subfigure[Intensity in 1999 - All Municipalities]
     {\includegraphics[width=0.49\textwidth]{figures/appendix/Figure_1b_inten1999_all.pdf}}
    \subfigure[Intensity in 2005 - All Municipalities]
     {\includegraphics[width=0.49\textwidth]{figures/appendix/Figure_1c_inten2005_all.pdf}}
      \subfigure[Intensity in 1997 - Highly Marginalized Municipalities]
     {\includegraphics[width=0.49\textwidth]{figures/appendix/Figure_1d_inten1997_mort.pdf}}
    \subfigure[Intensity in 1997 - All Municipalities]
     {\includegraphics[width=0.49\textwidth]{figures/appendix/Figure_1e_inten1997_mort_all.pdf}}
    
     
         \floatfoot{Notes: }
\end{figure}

\begin{figure}[H]
    \caption{The effects of interaction terms between Progresa in 1999 and years on elderly mortality +65 overall and by sex in marginalized municipalities, Mexico, 1991-2006. }
    \label{f:es_dd_sex_no_weight}

  \subfigure[Unweighted]
{\includegraphics[width=0.89\textwidth]{figures/appendix/Figure_2a_uw.pdf}}
  \subfigure[AAMR Weighted +SP]
{\includegraphics[width=0.89\textwidth]{figures/appendix/Figure_3_aamr_col4.pdf}}
    
  

    
         
         \floatfoot{Notes: The data analysis is restricted to only marginalized municipalities, 1991-2006. The plots present the interaction terms between Progresa intensity in 1999 and years. Year and municipality fixed effects are included. Standard errors are clustered at the municipality level. 1996 is a reference year and the regression models are not weighted.
Source: Authors’ calculations using data on deaths and population from the Instituto Nacional de Estadistica y Geografia and Mexican Census data.
}
\end{figure}



\begin{figure}[H]
    \caption{The effects of interaction terms between Progresa in 1999 and years on elderly mortality +65 overall and by sex in marginalized municipalities, Mexico, 1991-2006. }
    \label{f:es_dd_sex_no_weight}

  \subfigure[BR - EMR - All]
{\includegraphics[width=0.49\textwidth]{figures/appendix/Figure_3_emr65_BR.pdf}}
    \subfigure[HM - EMR - All]
{\includegraphics[width=0.49\textwidth]{figures/appendix/Figure_3_emr65_HighMarg.pdf}}


\subfigure[BR - EMR - Females]
{\includegraphics[width=0.49\textwidth]{figures/appendix/Figure_3_emr65f_BR.pdf}}
\subfigure[HM - EMR - Females]
     {\includegraphics[width=0.49\textwidth]{figures/appendix/Figure_3_emr65f_HighMarg.pdf}}

\subfigure[BR - EMR - Males]
{\includegraphics[width=0.49\textwidth]{figures/appendix/Figure_3_emr65m_BR.pdf}}
 \subfigure[HM - EMR - Males]
{\includegraphics[width=0.49\textwidth]{figures/appendix/Figure_3_emr65m_HighMarg.pdf}}
    
         \floatfoot{Notes: The data analysis is restricted to only marginalized municipalities, 1991-2006. The plots present the interaction terms between Progresa intensity in 1999 and years. Year and municipality fixed effects are included. Standard errors are clustered at the municipality level. 1996 is a reference year and the regression models are not weighted.
Source: Authors’ calculations using data on deaths and population from the Instituto Nacional de Estadistica y Geografia and Mexican Census data.
}
\end{figure}



% \begin{figure}[H]
%     \caption{The effects of interaction terms between Progresa in 1999 and years on elderly mortality +65 overall and by sex in marginalized municipalities, Mexico, 1991-2006. }
%     \label{f:es_dd_sex_no_weight}

%   \subfigure[BR - AAMR - All]
% {\includegraphics[width=0.49\textwidth]{figures/appendix/Figure_3_aamr65_BR.pdf}}
%     \subfigure[HM - AAMR - All]
% {\includegraphics[width=0.49\textwidth]{figures/appendix/Figure_3_aamr65_HighMarg.pdf}}


% \subfigure[BR - AAMR - Females]
% {\includegraphics[width=0.49\textwidth]{figures/appendix/Figure_3_aamr65f_BR.pdf}}
% \subfigure[HM - AAMR - Females]
%      {\includegraphics[width=0.49\textwidth]{figures/appendix/Figure_3_aamr65f_HighMarg.pdf}}

% \subfigure[BR - AAMR - Males]
% {\includegraphics[width=0.49\textwidth]{figures/appendix/Figure_3_aamr65m_BR.pdf}}
%  \subfigure[HM - AAMR - Males]
% {\includegraphics[width=0.49\textwidth]{figures/appendix/Figure_3_aamr65m_HighMarg.pdf}}
    
%          \floatfoot{Notes: The data analysis is restricted to only marginalized municipalities, 1991-2006. The plots present the interaction terms between Progresa intensity in 1999 and years. Year and municipality fixed effects are included. Standard errors are clustered at the municipality level. 1996 is a reference year and the regression models are not weighted.
% Source: Authors’ calculations using data on deaths and population from the Instituto Nacional de Estadistica y Geografia and Mexican Census data.
% }
% \end{figure}




\begin{figure}[H]
    \caption{The effects of interaction terms between Progresa in 1999 and years on elderly mortality +65 overall and by sex in marginalized municipalities, Mexico, 1991-2006. }
    \label{f:es_dd_sex_no_weight}

  \subfigure[Cancer]
{\includegraphics[width=0.49\textwidth]{figures/appendix/Figure_4_tb_cancer_Marg.pdf}}
\subfigure[Cardiovascular]
{\includegraphics[width=0.49\textwidth]{figures/appendix/Figure_4_tb_card_Marg.pdf}}

    \subfigure[Respiratory]
{\includegraphics[width=0.49\textwidth]{figures/appendix/Figure_4_tb_resp_Marg.pdf}}
    \subfigure[Infectious]
     {\includegraphics[width=0.49\textwidth]{figures/appendix/Figure_4_tb_infect_Marg.pdf}}
    

     
         \floatfoot{Notes: The data analysis is restricted to only marginalized municipalities, 1991-2006. The plots present the interaction terms between Progresa intensity in 1999 and years. Year and municipality fixed effects are included. Standard errors are clustered at the municipality level. 1996 is a reference year and the regression models are not weighted.
Source: Authors’ calculations using data on deaths and population from the Instituto Nacional de Estadistica y Geografia and Mexican Census data.
}
\end{figure}


\begin{figure}[H]
    \caption{The effects of interaction terms between Progresa in 1999 and years on elderly mortality +65 overall and by sex in marginalized municipalities, Mexico, 1991-2006. }
    \label{f:es_dd_sex_no_weight}



         \subfigure[Diabetes]
{\includegraphics[width=0.49\textwidth]{figures/appendix/Figure_4_tb_diab_Marg.pdf}}
    \subfigure[Nutri]
     {\includegraphics[width=0.49\textwidth]{figures/appendix/Figure_4_tb_nutri_Marg.pdf}}

              \subfigure[Ill-Defined]
{\includegraphics[width=0.49\textwidth]{figures/appendix/Figure_4_tb_illdef_Marg.pdf}}
    \subfigure[Accidents]
     {\includegraphics[width=0.49\textwidth]{figures/appendix/Figure_4_tb_accid_Marg.pdf}}

     

     
         \floatfoot{Notes: The data analysis is restricted to only marginalized municipalities, 1991-2006. The plots present the interaction terms between Progresa intensity in 1999 and years. Year and municipality fixed effects are included. Standard errors are clustered at the municipality level. 1996 is a reference year and the regression models are not weighted.
Source: Authors’ calculations using data on deaths and population from the Instituto Nacional de Estadistica y Geografia and Mexican Census data.
}
\end{figure}



\begin{figure}[H]
    \caption{The effects of interaction terms between Progresa in 1999 and years on elderly mortality +65 overall and by sex in marginalized municipalities, Mexico, 1991-2006. }
    \label{f:es_dd_sex_no_weight}

  \subfigure[Cancer]
{\includegraphics[width=0.49\textwidth]{figures/appendix/Figure_4_tb_cancer_BR.pdf}}
\subfigure[Cardiovascular]
{\includegraphics[width=0.49\textwidth]{figures/appendix/Figure_4_tb_card_BR.pdf}}

    \subfigure[Respiratory]
{\includegraphics[width=0.49\textwidth]{figures/appendix/Figure_4_tb_resp_BR.pdf}}
    \subfigure[Infectious]
     {\includegraphics[width=0.49\textwidth]{figures/appendix/Figure_4_tb_infect_BR.pdf}}


     

     
         \floatfoot{Notes: The data analysis is restricted to only marginalized municipalities, 1991-2006. The plots present the interaction terms between Progresa intensity in 1999 and years. Year and municipality fixed effects are included. Standard errors are clustered at the municipality level. 1996 is a reference year and the regression models are not weighted.
Source: Authors’ calculations using data on deaths and population from the Instituto Nacional de Estadistica y Geografia and Mexican Census data.
}
\end{figure}


\begin{figure}[H]
    \caption{The effects of interaction terms between Progresa in 1999 and years on elderly mortality +65 overall and by sex in marginalized municipalities, Mexico, 1991-2006. }
    \label{f:es_dd_sex_no_weight}

         \subfigure[Diabetes]
{\includegraphics[width=0.49\textwidth]{figures/appendix/Figure_4_tb_diab_BR.pdf}}
    \subfigure[Nutri]
     {\includegraphics[width=0.49\textwidth]{figures/appendix/Figure_4_tb_nutri_BR.pdf}}

              \subfigure[Ill-Defined]
{\includegraphics[width=0.49\textwidth]{figures/appendix/Figure_4_tb_illdef_BR.pdf}}
    \subfigure[Accidents]
     {\includegraphics[width=0.49\textwidth]{figures/appendix/Figure_4_tb_accid_BR.pdf}}

     

     
         \floatfoot{Notes: The data analysis is restricted to only marginalized municipalities, 1991-2006. The plots present the interaction terms between Progresa intensity in 1999 and years. Year and municipality fixed effects are included. Standard errors are clustered at the municipality level. 1996 is a reference year and the regression models are not weighted.
Source: Authors’ calculations using data on deaths and population from the Instituto Nacional de Estadistica y Geografia and Mexican Census data.
}
\end{figure}



\newcommand{\mectable}[4]{%
  \begin{table}[H]
    \centering
    \caption{#1}
    \label{#2}
    \resizebox{\textwidth}{!}{\input{#3}}
    {\footnotesize\raggedright\noindent Notes: #4\par}
  \end{table}
  \clearpage
}

\begin{table}[H]
\RawCaption%
  {\caption{\deleted{Barham and Rowberry’s results and our replication: Impact of Progresa on elderly mortality 65+, municipalities incorporated in 1998 or 1999}\added{Replication of Barham and Rowberry (2013): Effect of Progresa on Elder Mortality Rate (EMR 65+), Municipalities Incorporated in 1998--1999}}
   \label{t:did_robust_br}}
\floatbox[{\captop}]{table}[\FBwidth]
  {}
  {\scalebox{0.9}{
      \begin{tabular}{lccc} \hline \hline
& \multicolumn{1}{c}{Pooled} & \multicolumn{1}{c}{Males} & \multicolumn{1}{c}{Females} \\ 
\cmidrule(lr){2-2}\cmidrule(lr){3-3}\cmidrule(lr){4-4}  
& \multicolumn{1}{c}{(1)} & \multicolumn{1}{c}{(2)} & \multicolumn{1}{c}{(3)} \\ \toprule 
\underline{\textit{Panel A: BR (2013)}}  \\ 
\textit{2-yr lagged Intensity} & -6.37*** & -6.42*** & -6.46*** \\ 
 & (1.04) & (1.42) & (1.31) \\ 
  & & & \\ 
Mean 1996 & 47.5 & 49.3 & 46.0 \\ 
Obs & 21,571 & 21,571 & 21,571 \\ 
No. Mun & 1,961 & 1,961 & 1,961 \\ 
  & & & \\ 
\underline{\textit{Panel B: Replication (Unweighted)}}  \\ 
\textit{2-yr lagged Intensity} &       -7.061*** &       -6.889*** &       -7.426*** \\ 
 & (       1.256) & (       1.711) & (       1.559) \\ 
  & & & \\ 
\underline{\textit{Panel C: Replication (Weighted)}}  \\ 
\textit{2-yr lagged Intensity} &       -3.740*** &       -3.322*** &       -4.210*** \\ 
 & (       0.786) & (       0.935) & (       1.002) \\ 
  & & & \\ 
\underline{\textit{Panel D: 1-yr Lag (Weighted)}}  \\ 
\textit{1-yr lagged Intensity} &       -3.457*** &       -2.681** &       -4.255*** \\ 
 & (       0.890) & (       1.072) & (       1.043) \\ 
  & & & \\ 
\underline{\textit{Panel E: 3-yr Lag (Weighted)}}  \\ 
\textit{3-yr lagged Intensity} &       -2.071*** &       -2.357** &       -1.830* \\ 
 & (       0.772) & (       0.957) & (       0.960) \\ 
  & & & \\ 
Mean 1996 &        48.16 &        48.87 &        47.49 \\ 
Obs &       15,642 &       15,642 &       15,642 \\ 
No. Mun &        1,422 &        1,422 &        1,422 \\ 
  & & & \\ 
Year FE & Y & Y & Y \\ 
Mun FE & Y & Y & Y \\ 
Cluster SE: Mun & Y & Y & Y \\ 
\bottomrule
\end{tabular}
    }
    \vspace{-3mm}
    \floatfoot{\deleted{Notes: The data analysis is restricted to municipalities incorporated into Progresa in 1998 or 1999. Progresa intensity data are pooled from 1990 to 2000. All regression models include year fixed effects and municipality fixed effects. Standard errors are clustered at the municipality level. The regression model is weighted by the number of population aged +65 for both sexes in (1), only male in (2), and only female in (3).\\ ***p$<0.01$, **p$<0.05$, *p$<0.1$.}\added{Notes: This table replicates the main specification of Barham and Rowberry (2013). The sample is restricted to municipalities incorporated into Progresa in 1998 or 1999. The outcome is the Elder Mortality Rate (EMR) for adults aged 65 and older per 1,000 inhabitants. Progresa intensity is measured using data pooled from 1990 to 2000. All regressions include municipality and year fixed effects. The regression model is weighted by the population aged 65 and older, shown separately for the pooled (column 1), male (column 2), and female (column 3) populations. Standard errors are clustered at the municipality level and shown in parentheses. ***p$<0.01$, **p$<0.05$, *p$<0.1$.}}
  }
\end{table}
\newpage




\begin{table}[H]
\RawCaption%
  {\caption{\deleted{Barham and Rowberry Extended Results Replication}\added{Sensitivity Analysis: Barham and Rowberry (2013) Replication with Alternative Lag Structures and Sample Restrictions}}
   \label{t:did_robust_br2}}
\floatbox[{\captop}]{table}[\FBwidth]
  {}
  {\scalebox{0.9}{
      \input{tables/appendix/T_BR_robustness_emr65}
    }
    \vspace{-3mm}
    \floatfoot{\deleted{Notes: Progresa intensity data are pooled from 1990 to 2000. All regression models include year fixed effects and municipality fixed effects. Standard errors are clustered at the municipality level. The regression model is weighted by the number of population aged +65. ***p$<0.01$, **p$<0.05$, *p$<0.1$.}\added{Notes: This table reports sensitivity analyses extending the replication of Barham and Rowberry (2013), varying the lag structure of Progresa intensity, weighting, and sample restrictions. The outcome is the Elder Mortality Rate (EMR) for adults aged 65 and older per 1,000 inhabitants. Row (A) reproduces Barham and Rowberry's (2013) original results; row (B) reports our replication using the same specification. Subsequent rows vary the lag structure (1-year and 3-year lags), weighting, and sample. Progresa intensity data are pooled from 1990 to 2000. All regressions include municipality and year fixed effects. Standard errors are clustered at the municipality level. ***p$<0.01$, **p$<0.05$, *p$<0.1$.}}
  }
\end{table}

% \begin{table}[H]
% \RawCaption%
%   {\caption{Barham and Rowberry Extended Results Replication}
%    \label{t:did_robust}}
% \floatbox[{\captop}]{table}[\FBwidth]
%   {}
%   {\scalebox{0.9}{
%       \input{tables/appendix/T_BR_robustness_aamr65}
%     }
%     \vspace{-3mm}
%     \floatfoot{Notes: Progresa intensity data are pooled from 1990 to 2000. All regression models include year fixed effects and municipality fixed effects. Standard errors are clustered at the municipality level. The regression model is weighted by the number of population aged +65. ***p$<0.01$, **p$<0.05$, *p$<0.1$.}
%   }
% \end{table}





% \begin{table}[H]
% \RawCaption%
%   {\caption{Functional Forms}
%    \label{t:did_robust}}
% \floatbox[{\captop}]{table}[\FBwidth]
%   {}
%   {\scalebox{0.9}{
%       \begin{tabular}{lccc} \hline \hline
& Levels & Log & Poisson \\ 
\cmidrule(lr){2-2}\cmidrule(lr){3-3}\cmidrule(lr){4-4}
& \multicolumn{1}{c}{(1)} & \multicolumn{1}{c}{(2)} & \multicolumn{1}{c}{(3)} \\ \toprule
\underline{\textit{Panel A: Pooled}}  \\ 
\textit{Intensity 1999 x post (1997-2006)} &       -0.777 &       -0.020 &       -0.031 \\ 
  & (       1.464) & (       0.042) & (       0.058) \\ 
   & & & \\ 
\textit{Intensity 2005 x post (1997-2006)} &       -2.787 &       -0.052 &       -0.128* \\ 
 & (       1.806) & (       0.052) & (       0.072) \\ 
  & & & \\ 
Mean (1991-1996)  &        46.47 &         3.75 &        23.86 \\ 
  & & &  \\ 
\underline{\textit{Panel B: Males}}  \\ 
\textit{Intensity 1999 x post (1997-2006)} &       -2.288 &       -0.029 &       -0.104* \\ 
  & (       1.732) & (       0.050) & (       0.062) \\ 
  & & & \\ 
\textit{Intensity 2005 x post (1997-2006)}  &        0.214 &       -0.020 &       -0.024 \\ 
 & (       2.000) & (       0.057) & (       0.076) \\ 
 & & & \\ 
Mean (1991-1996) &        47.26 &         3.78 &        12.25 \\ 
  & & &  \\ 
\underline{\textit{Panel C: Females}}  \\ 
\textit{Intensity 1999 x post (1997-2006)}  &        0.592 &        0.025 &        0.055 \\ 
  & (       1.824) & (       0.055) & (       0.074) \\ 
  & & & \\ 
\textit{Intensity 2005 x post (1997-2006)} &       -5.637** &       -0.173** &       -0.223*** \\ 
 & (       2.205) & (       0.067) & (       0.076) \\ 
   & & & \\ 
Mean (1991-1996)  &        46.11 &         3.73 &        11.61 \\ 
  & & &  \\ 
Obs &       17,904 &       17,223 &       17,904 \\ 
No. Mun &        1,119 &        1,119 &        1,119 \\ 
  & & & \\ 
Year FE & Y & Y & Y \\ 
Mun FE  & Y & Y & Y \\ 
Weights & Y & Y & Y \\ 
Cluster SE: Mun & Y & Y & Y \\ 
\bottomrule
\end{tabular}
%     }
%     \vspace{-3mm}
%     \floatfoot{Notes: Progresa intensity data are pooled from 1990 to 2000. All regression models include year fixed effects and municipality fixed effects. Standard errors are clustered at the municipality level. The regression model is weighted by the number of population aged +65. ***p$<0.01$, **p$<0.05$, *p$<0.1$.}
%   }
% \end{table}
% \newpage






\begin{table}[H]
\RawCaption%
  {\caption{\deleted{AAMR}\added{Effects of Progresa Enrollment Intensity on Age-Adjusted Mortality Rate (AAMR 65+) by Sex, Marginalized Municipalities, 1991--2006}}
   \label{t:did_robust_aamr}}
\floatbox[{\captop}]{table}[\FBwidth]
  {}
  {\scalebox{0.9}{
      \input{tables/appendix/AT4_aamr_mortality}
    }
    \vspace{-3mm}
    \floatfoot{\deleted{Notes: Progresa intensity data are pooled from 1990 to 2000. All regression models include year fixed effects and municipality fixed effects. Standard errors are clustered at the municipality level. The regression model is weighted by the number of population aged +65. ***p$<0.01$, **p$<0.05$, *p$<0.1$.}\added{Notes: This table reports difference-in-differences estimates of the effect of Progresa enrollment intensity on the Age-Adjusted Mortality Rate (AAMR) for adults aged 65 and older per 1,000 inhabitants, separately for the pooled, male, and female populations. Columns (1) and (2) report unweighted estimates without and with Seguro Popular intensity; columns (3) and (4) report estimates weighted by the population aged 65 and older without and with Seguro Popular intensity. The post-dummy equals 1 for years 1997--2006. The sample is restricted to highly marginalized municipalities ($\textit{gm\_mun\_1990} = 4$ or $5$). All regressions include municipality and year fixed effects. Standard errors are clustered at the municipality level. ***p$<0.01$, **p$<0.05$, *p$<0.1$.}}
  }
\end{table}

\begin{landscape}
\begin{table}[H]
\RawCaption%
  {\caption{\deleted{AAMR}\added{Effects of Progresa Enrollment Intensity on Cause-Specific Elder Mortality Rate (EMR 65+), Highly Marginalized Municipalities, 1991--2006}}
   \label{t:did_robust_cod}}
\floatbox[{\captop}]{table}[\FBwidth]
  {}
  {\scalebox{0.7}{
      \input{tables/appendix/AT5_cod_mortality}
    }
    \vspace{-3mm}
    \floatfoot{\deleted{Notes: Progresa intensity data are pooled from 1990 to 2000. All regression models include year fixed effects and municipality fixed effects. Standard errors are clustered at the municipality level. The regression model is weighted by the number of population aged +65. ***p$<0.01$, **p$<0.05$, *p$<0.1$.}\added{Notes: This table reports difference-in-differences estimates of the effect of Progresa enrollment intensity on cause-specific Elder Mortality Rate (EMR) for adults aged 65 and older per 1,000 inhabitants. Each column corresponds to a specific cause of death. The post-dummy equals 1 for years 1997--2006. The sample is restricted to highly marginalized municipalities ($\textit{gm\_mun\_1990} = 4$ or $5$). All regressions include municipality and year fixed effects. Standard errors are clustered at the municipality level. ***p$<0.01$, **p$<0.05$, *p$<0.1$.}}
  }
\end{table}
\end{landscape}

\begin{landscape}
\begin{table}[H]
\RawCaption%
  {\caption{\deleted{AAMR}\added{Effects of Progresa Enrollment Intensity on Cause-Specific Elder Mortality Rate (EMR 65+), Barham-Rowberry Sample, 1992--2002}}
   \label{t:did_robust_cod_br}}
\floatbox[{\captop}]{table}[\FBwidth]
  {}
  {\scalebox{0.9}{
      \begin{tabular}{lcccccccccc} \hline \hline
& Cancer & Diab. & IllDef & Resp. & Card. & Infect. & Nutri. & Accid. & Other \\ 
\cmidrule(lr){2-2}\cmidrule(lr){3-3}\cmidrule(lr){4-4}\cmidrule(lr){5-5}\cmidrule(lr){6-6}\cmidrule(lr){7-7}\cmidrule(lr){8-8}\cmidrule(lr){9-9}\cmidrule(lr){10-10}
& \multicolumn{1}{c}{(1)} & \multicolumn{1}{c}{(2)} & \multicolumn{1}{c}{(3)} & \multicolumn{1}{c}{(4)} & \multicolumn{1}{c}{(5)} & \multicolumn{1}{c}{(6)} & \multicolumn{1}{c}{(7)} & \multicolumn{1}{c}{(8)} & \multicolumn{1}{c}{(9)} \\ \toprule
\underline{\textit{Panel A: Pooled}}  \\ 
\textit{Intensity 1999 x post} &       -0.816*** &       -0.924*** &       -0.614*** &        0.149 &        0.792* &       -1.670*** &       -0.279 &        0.019 &        1.681*** \\ 
 & (       0.188) & (       0.183) & (       0.134) & (       0.281) & (       0.470) & (       0.210) & (       0.248) & (       0.048) & (       0.498) \\ 
\textit{RW p-value} & [ 0.000] & [ 0.000] & [ 0.000] & [ 1.000] & [ 0.370] & [ 0.000] & [ 0.784] & [ 1.000] & [ 0.004] \\ 
  & & & & & & & & & \\ 
Mean (pre-1997) &         2.48 &         2.58 &         0.91 &         4.86 &        15.84 &         2.56 &         3.54 &         0.27 &        14.39 \\ 
Obs &       15,642 &       15,642 &       15,642 &       15,642 &       15,642 &       15,642 &       15,642 &       15,642 &       15,642 \\ 
  & & & & & & & & & \\ 
\underline{\textit{Panel B: Females}}  \\ 
\textit{Intensity 1999 x post} &       -0.600*** &       -0.758*** &       -0.687*** &       -0.126 &        0.821 &       -1.792*** &       -0.447 &       -0.035 &        1.279** \\ 
 & (       0.215) & (       0.271) & (       0.170) & (       0.350) & (       0.577) & (       0.227) & (       0.308) & (       0.043) & (       0.579) \\ 
\textit{RW p-value} & [ 0.037] & [ 0.037] & [ 0.000] & [ 0.845] & [ 0.591] & [ 0.000] & [ 0.591] & [ 0.845] & [ 0.137] \\ 
  & & & & & & & & & \\ 
Mean (pre-1997) &         2.13 &         3.15 &         0.94 &         4.43 &        16.35 &         2.46 &         3.86 &         0.12 &        13.25 \\ 
Obs &       15,642 &       15,642 &       15,642 &       15,642 &       15,642 &       15,642 &       15,642 &       15,642 &       15,642 \\ 
  & & & & & & & & & \\ 
\underline{\textit{Panel C: Males}}  \\ 
\textit{Intensity 1999 x post} &       -1.074*** &       -1.086*** &       -0.559*** &        0.450 &        0.702 &       -1.532*** &       -0.140 &        0.075 &        2.163*** \\ 
 & (       0.262) & (       0.212) & (       0.159) & (       0.334) & (       0.565) & (       0.261) & (       0.263) & (       0.083) & (       0.641) \\ 
\textit{RW p-value} & [ 0.000] & [ 0.000] & [ 0.003] & [ 0.714] & [ 0.714] & [ 0.000] & [ 0.734] & [ 0.734] & [ 0.004] \\ 
  & & & & & & & & & \\ 
Mean (pre-1997) &         2.85 &         2.04 &         0.89 &         5.28 &        15.38 &         2.69 &         3.21 &         0.42 &        15.73 \\ 
Obs &       15,642 &       15,642 &       15,642 &       15,642 &       15,642 &       15,642 &       15,642 &       15,642 &       15,642 \\ 
  & & & & & & & & & \\ 
\bottomrule
\end{tabular}
    }
    \vspace{-3mm}
    \floatfoot{\deleted{Notes: Progresa intensity data are pooled from 1990 to 2000. All regression models include year fixed effects and municipality fixed effects. Standard errors are clustered at the municipality level. The regression model is weighted by the number of population aged +65. ***p$<0.01$, **p$<0.05$, *p$<0.1$.}\added{Notes: This table reports difference-in-differences estimates of the effect of Progresa enrollment intensity on cause-specific Elder Mortality Rate (EMR) for adults aged 65 and older per 1,000 inhabitants, using the Barham-Rowberry sample (municipalities incorporated into Progresa in 1998 or 1999). Each column corresponds to a specific cause of death. All regressions include municipality and year fixed effects. Standard errors are clustered at the municipality level. ***p$<0.01$, **p$<0.05$, *p$<0.1$.}}
  }
\end{table}
\end{landscape}


\begin{table}[htbp]
\centering
\caption{\deleted{Progresa Impacts on Household Expenditures: Heterogeneity by Elderly Presence (Eligible Households)}\added{Short-Run Effects of Progresa on Household Expenditure by Elderly Presence, Eligible Households}}
\label{tab:expenditures_elderly}
\floatbox[{\captop}]{table}[\FBwidth]
  {}
  {\scalebox{0.9}{
      \begin{tabular}{lcccc}
\toprule
 & (1) & (2) & (3) & (4) \\
 & ln Food & ln Medicine & \% Food & \% Medicine \\
\midrule
Treat $\times$ 1999 (No Elderly)   & 0.121$^{***}$ & 0.039 & $-$0.014 & 0.000 \\
                                    & (0.026) & (0.062) & (0.014) & (0.005) \\[4pt]
Diff: Elderly HH $\times$ 1999     & 0.007 & $-$0.012 & 0.021 & 0.004 \\
                                    & (0.035) & (0.115) & (0.017) & (0.011) \\[4pt]
Observations                        & 22,524 & 22,763 & 22,478 & 22,489 \\
Control Mean (1997)                 & 4.668 & 0.744 & 0.757 & 0.053 \\
\bottomrule
\end{tabular}



    }
    \vspace{-3mm}
    \floatfoot{\deleted{Notes: . ***p$<0.01$, **p$<0.05$, *p$<0.1$.}\added{Notes: This table reports difference-in-differences estimates of the effect of Progresa enrollment intensity on household expenditure outcomes, stratified by the presence of elderly members in the household. The sample is restricted to eligible households. All regressions include municipality and year fixed effects. Standard errors are clustered at the municipality level and shown in parentheses. ***p$<0.01$, **p$<0.05$, *p$<0.1$.}}
  }
\end{table}

\end{document}
